% Mathématiques 5e — cours V3
% Nombres relatifs

\section{Objectifs}

\begin{itemize}
\item Définir les nombres relatifs.
\item Distinguer nombres positifs, négatifs, strictement positifs et strictement négatifs.
\item Définir l'opposé et la valeur absolue.
\item Utiliser des nombres relatifs pour modéliser des températures, altitudes, écarts ou chronologies.
\item Lire et placer un nombre relatif sur une droite graduée.
\item Comparer et ranger des nombres décimaux relatifs.
\item Additionner deux ou plusieurs nombres décimaux relatifs.
\item Soustraire deux nombres décimaux relatifs.
\item Comprendre le rôle des parenthèses dans les écritures comportant des signes.
\item Enchainer additions et soustractions.
\item Résoudre des problèmes mobilisant des nombres relatifs.
\end{itemize}

\section{1. Pourquoi introduire des nombres négatifs ?}

Les nombres déjà connus ne permettent pas toujours de décrire toutes les situations.

Une température peut être inférieure à :
\[
0^\circ\text{C}.
\]

Une altitude peut être située sous le niveau de la mer.

Un compte peut présenter un solde négatif.

Pour représenter ces valeurs, on utilise les nombres relatifs.

\coursefigure{/images/mathematiques/5eme/relatifs-contexte.svg}{Trois contextes utilisant des valeurs positives ou négatives autour d'une référence zéro.}{Les nombres relatifs permettent de décrire des grandeurs situées de part et d'autre d'une valeur de référence.}

\section{2. Définition}

\begin{definition}
Un \textbf{nombre relatif} est un nombre muni d'un signe positif ou négatif, ou égal à zéro.
\end{definition}

Exemples :
\[
+7,\qquad -3,\qquad +2{,}5,\qquad -8{,}4,\qquad 0.
\]

On peut omettre le signe $+$ devant un nombre positif :
\[
+7=7.
\]

Le signe $-$ fait partie du nombre.

\section{3. Positif, négatif, strictement positif}

Un nombre est positif s'il est supérieur ou égal à zéro :
\[
x\ge0.
\]

Un nombre est strictement positif s'il est supérieur à zéro :
\[
x>0.
\]

Un nombre est négatif s'il est inférieur ou égal à zéro :
\[
x\le0.
\]

Un nombre est strictement négatif s'il est inférieur à zéro :
\[
x<0.
\]

\begin{attention}
Le nombre :
\[
0
\]
est à la fois positif et négatif au sens large, mais il n'est ni strictement positif ni strictement négatif.
\end{attention}

\section{4. Droite graduée}

Les nombres relatifs sont placés sur une droite graduée.

\coursefigure{/images/mathematiques/5eme/relatifs-droite.svg}{Droite graduée de moins quatre à quatre avec moins deux et deux mis en évidence.}{Sur une droite graduée, les nombres augmentent de gauche à droite ; deux nombres opposés sont symétriques par rapport à zéro.}

À droite de zéro se trouvent les nombres positifs.

À gauche de zéro se trouvent les nombres négatifs.

Sur une droite graduée :
\[
-4<-1<0<2<5.
\]

\section{5. Lire une abscisse}

L'abscisse d'un point est le nombre qui repère sa position sur la droite.

Si un point $A$ est placé sur la graduation :
\[
-3,
\]
on écrit :
\[
A(-3).
\]

Si $B$ est placé sur :
\[
2{,}5,
\]
on écrit :
\[
B(2{,}5).
\]

\section{6. Placer un nombre relatif}

Pour placer :
\[
-2{,}5,
\]
on part de zéro puis on se déplace de :
\[
2{,}5
\]
unités vers la gauche.

Pour placer :
\[
+3{,}2,
\]
on se déplace de :
\[
3{,}2
\]
unités vers la droite.

Il faut toujours repérer l'unité de graduation avant de placer un point.

\section{7. Opposé}

\begin{definition}
Deux nombres sont \textbf{opposés} lorsqu'ils ont la même distance à zéro mais sont situés de part et d'autre de zéro.
\end{definition}

L'opposé de :
\[
5
\]
est :
\[
-5.
\]

L'opposé de :
\[
-3{,}2
\]
est :
\[
3{,}2.
\]

L'opposé de zéro est zéro.

\begin{propriete}
La somme de deux nombres opposés vaut :
\[
0.
\]
\end{propriete}

Ainsi :
\[
7+(-7)=0.
\]

\section{8. Valeur absolue}

\begin{definition}
La valeur absolue d'un nombre relatif est sa distance à zéro.
\end{definition}

On note :
\[
|x|.
\]

Par exemple :
\[
|-5|=5,
\]
et :
\[
|3{,}2|=3{,}2.
\]

La valeur absolue est toujours positive ou nulle :
\[
|x|\ge0.
\]

\section{9. Comparer deux nombres relatifs}

Sur une droite graduée, le plus grand nombre est celui situé le plus à droite.

Par exemple :
\[
-7<-2.
\]

Même si :
\[
7>2,
\]
on a pour les nombres négatifs :
\[
-7<-2.
\]

Le nombre le plus éloigné de zéro du côté négatif est le plus petit.

\section{10. Comparer un négatif et un positif}

Tout nombre strictement négatif est inférieur à tout nombre strictement positif.

Par exemple :
\[
-0{,}1<2.
\]

De même :
\[
-100<0{,}0001.
\]

Le signe suffit ici pour conclure.

\section{11. Ranger une liste}

Ranger dans l'ordre croissant :
\[
3{,}2,\quad -4,\quad 0,\quad -1{,}5,\quad 2.
\]

On obtient :
\[
\boxed{-4<-1{,}5<0<2<3{,}2}.
\]

\begin{methode}
Pour ranger des nombres relatifs :
\begin{enumerate}
\item placer mentalement ou réellement les valeurs sur une droite ;
\item traiter les nombres négatifs ;
\item placer zéro ;
\item traiter les nombres positifs.
\end{enumerate}
\end{methode}

\section{12. Additionner deux nombres de même signe}

Lorsque deux nombres sont positifs, on additionne normalement :
\[
(+4)+(+3)=+7.
\]

Lorsque deux nombres sont négatifs, on additionne leurs distances à zéro et on conserve le signe négatif :
\[
(-4)+(-3)=-7.
\]

\begin{propriete}
Pour additionner deux nombres de même signe :
\begin{enumerate}
\item on additionne les valeurs absolues ;
\item on conserve le signe commun.
\end{enumerate}
\end{propriete}

\section{13. Additionner deux nombres de signes contraires}

Considérons :
\[
(+8)+(-5).
\]

Les nombres ont des signes contraires.

On calcule la différence des valeurs absolues :
\[
8-5=3.
\]

Le nombre de plus grande valeur absolue est positif.

Donc :
\[
(+8)+(-5)=+3.
\]

Autre exemple :
\[
(-9)+(+4).
\]

On calcule :
\[
9-4=5.
\]

Le nombre de plus grande valeur absolue est négatif.

Donc :
\[
(-9)+(+4)=-5.
\]

\section{14. Modèle du déplacement}

L'addition peut être interprétée comme un déplacement sur la droite graduée.

À partir de :
\[
-2,
\]
ajouter :
\[
+5
\]
revient à se déplacer de cinq unités vers la droite.

On arrive à :
\[
3.
\]

Donc :
\[
-2+5=3.
\]

Ajouter un nombre négatif correspond à se déplacer vers la gauche.

\section{15. Additionner plusieurs nombres relatifs}

Considérons :
\[
-3+7-2+5.
\]

On peut enchainer de gauche à droite :
\[
-3+7=4,
\]
puis :
\[
4-2=2,
\]
puis :
\[
2+5=7.
\]

Donc :
\[
\boxed{-3+7-2+5=7}.
\]

On peut parfois regrouper intelligemment :
\[
-3-2+7+5.
\]

Alors :
\[
-5+12=7.
\]

\section{16. Soustraire un nombre relatif}

Soustraire un nombre revient à ajouter son opposé.

\begin{propriete}
Pour tous nombres $a$ et $b$ :
\[
a-b=a+(-b).
\]
\end{propriete}

Par exemple :
\[
5-(-3)
=
5+3
=
8.
\]

Et :
\[
-4-(+2)
=
-4+(-2)
=
-6.
\]

\section{17. Pourquoi les parenthèses sont parfois indispensables}

Dans :
\[
5-(-3),
\]
les parenthèses distinguent clairement :
\begin{itemize}
\item le signe de l'opération de soustraction ;
\item le signe du nombre $-3$.
\end{itemize}

Sans parenthèses, l'écriture :
\[
5--3
\]
est difficile à lire et n'est pas l'écriture usuelle attendue.

\begin{attention}
Les parenthèses ne sont pas décoratives. Elles peuvent être indispensables pour comprendre quel signe appartient au nombre et quel signe indique une opération.
\end{attention}

\section{18. Simplifier une écriture}

On peut simplifier certaines écritures.

Par exemple :
\[
7+(+3)=7+3.
\]

De même :
\[
7+(-3)=7-3.
\]

Et :
\[
7-(-3)=7+3.
\]

Ces transformations doivent conserver exactement le même nombre.

\section{19. Exemple développé : température}

À $6$ h, la température est :
\[
-4^\circ\text{C}.
\]

Elle augmente ensuite de :
\[
7^\circ\text{C}.
\]

La nouvelle température vaut :
\[
-4+7=3.
\]

Donc :
\[
\boxed{3^\circ\text{C}}.
\]

Puis elle baisse de :
\[
5^\circ\text{C}.
\]

On calcule :
\[
3-5=-2.
\]

La température devient :
\[
\boxed{-2^\circ\text{C}}.
\]

\section{20. Exemple développé : altitude}

Un plongeur se situe à :
\[
-18\ \text{m}
\]
par rapport au niveau de la mer.

Il remonte de :
\[
7\ \text{m}.
\]

Son altitude devient :
\[
-18+7=-11.
\]

Il se trouve donc à :
\[
\boxed{-11\ \text{m}}.
\]

\section{21. Écart entre deux valeurs}

La différence entre deux températures peut être calculée par une soustraction.

Entre :
\[
-3^\circ\text{C}
\]
et :
\[
5^\circ\text{C},
\]
l'écart vaut :
\[
5-(-3)=8.
\]

La différence de température est :
\[
\boxed{8^\circ\text{C}}.
\]

\section{22. Valeur absolue et distance}

La distance entre zéro et :
\[
-6
\]
vaut :
\[
|-6|=6.
\]

La valeur absolue ne conserve donc pas le signe.

Elle mesure uniquement une distance à zéro.

\section{23. Limite du cours de 5e}

\begin{attention}
En 5e, l'objectif porte sur l'addition et la soustraction des nombres décimaux relatifs.

Les règles générales de signe pour la multiplication et la division de deux nombres relatifs sont étudiées en 4e.
\end{attention}

On ne doit donc pas anticiper inutilement :
\[
(-3)\times(-5)
\]
comme automatisme du chapitre de 5e.

\section{24. Méthode générale pour une somme}

\begin{methode}
Pour additionner deux nombres relatifs :
\begin{enumerate}
\item repérer leurs signes ;
\item s'ils ont le même signe, additionner les valeurs absolues et conserver ce signe ;
\item s'ils ont des signes différents, calculer la différence des valeurs absolues ;
\item donner au résultat le signe du nombre de plus grande valeur absolue ;
\item vérifier le résultat sur une droite graduée si nécessaire.
\end{enumerate}
\end{methode}

\section{25. Méthode générale pour une soustraction}

\begin{methode}
Pour calculer :
\[
a-b,
\]
on transforme la soustraction en addition :
\[
a-b=a+(-b).
\]

On applique ensuite la méthode d'addition des nombres relatifs.
\end{methode}

\section{26. Erreurs fréquentes}

\begin{itemize}
\item Croire que $-8$ est plus grand que $-3$ parce que $8>3$.
\item Confondre opposé et valeur absolue.
\item Oublier que $0$ n'est ni strictement positif ni strictement négatif.
\item Additionner les valeurs absolues lorsque les signes sont différents.
\item Transformer incorrectement une soustraction.
\item Supprimer des parenthèses sans tenir compte des signes.
\item Utiliser prématurément les règles de multiplication des nombres relatifs.
\end{itemize}

\section{À retenir}

\begin{aretenir}
Les nombres relatifs permettent de décrire des valeurs situées de part et d'autre de zéro.

Deux nombres opposés ont la même valeur absolue et des signes contraires.

La valeur absolue mesure la distance à zéro :
\[
\boxed{|x|\ge0}.
\]

Pour une soustraction :
\[
\boxed{a-b=a+(-b)}.
\]

Sur une droite graduée, le plus grand nombre est toujours celui situé le plus à droite.

En 5e, on automatise surtout :
\[
\boxed{\text{comparaison, addition et soustraction des nombres relatifs}}.
\]
\end{aretenir}
